\documentclass[10pt]{article}
\usepackage[T1]{fontenc}
\usepackage[utf8]{inputenc}
\usepackage[french]{babel}
\usepackage{amsmath,amssymb,mathrsfs,mathpazo}
\usepackage{geometry}
\geometry{a4paper,margin=2cm}
\usepackage{graphicx,array,tabularx,tabularray,colortbl,multirow,multicol,booktabs,enumitem,pifont,ulem,fancybox,diagbox,cancel}
\usepackage{pstricks,pstricks-add,pst-plot,pst-tree,pst-node,pst-text,pst-func,pst-math,pst-bspline,pst-3dplot,pst-eucl,pst-solides3d}
\let\pscalculate\undefined
\usepackage{pst-bezier}
\usepackage{tikz,pgfplots,tkz-tab}
\usetikzlibrary{arrows,arrows.meta,patterns,calc,positioning,trees,decorations.markings,shapes.geometric,shapes.misc,shapes.arrows,intersections,decorations.pathmorphing,decorations.pathreplacing,backgrounds,shadows,babel}
\pgfplotsset{compat=1.18}
\usepgfplotslibrary{fillbetween}
\usepackage[np]{numprint}
\usepackage{listings,makecell,pas-tableur,twoopt}
\let\barre\undefined
\usepackage[pstricks]{tablvar}
\definecolor{darkgreen}{RGB}{0,100,0}
\providecommand{\red}{\color{red}}
\providecommand{\blue}{\color{blue}}
\providecommand{\green}{\color{green}}
\providecommand{\black}{\color{black}}
\providecommand{\white}{\color{white}}
\providecommand{\gray}{\color{gray}}
\providecommand{\cyan}{\color{cyan}}
\providecommand{\magenta}{\color{magenta}}
\providecommand{\yellow}{\color{yellow}}
\providecommand{\violet}{\color{violet}}
\providecommand{\orange}{\color{orange}}
\providecommand{\euro}{\texteuro}
\providecommand{\EUR}{\texteuro}
\providecommand{\ieme}{\textsuperscript{e}}
\providecommand{\degres}{\ensuremath{^\circ}}
\providecommand{\degree}{\ensuremath{^\circ}}
\providecommand{\vv}[1]{\overrightarrow{#1}}
\providecommand{\Coord}[1]{#1}
\providecommand{\decosix}{}
\providecommand{\starredbullet}{\ensuremath{\bullet}}
\pagestyle{empty}
\setlength{\parindent}{0pt}
\usepackage{adjustbox}\DeclareUnicodeCharacter{2236}{:}\DeclareUnicodeCharacter{2190}{\ensuremath{\leftarrow}}
\providecommand{\R}{\mathbb{R}}
\providecommand{\N}{\mathbb{N}}
\providecommand{\Z}{\mathbb{Z}}
\providecommand{\Q}{\mathbb{Q}}
\providecommand{\D}{\mathbb{D}}
\providecommand{\C}{\mathbb{C}}
\providecommand{\e}{{\mathrm{e}}}
\providecommand{\vect}[1]{\overrightarrow{#1}}
\providecommand{\vectt}[1]{\overrightarrow{\mathrm{#1}}}
\providecommand{\vvt}[1]{\overrightarrow{\mathrm{#1}}}
\providecommand{\np}[2][]{{#2~#1}}
\providecommand{\vv}[1]{{\overrightarrow{#1}}}
\renewcommand{\d}{{\mathrm{d}}}
\providecommand{\degres}{{^\circ}}
\providecommand{\degre}{{^\circ}}
\providecommand{\blue}{\color{blue}}
\providecommand{\red}{\color{red}}
\providecommand{\enskip}{\quad}
\providecommand{\llbracket}{\lbrack\!\lbrack}
\providecommand{\rrbracket}{\rbrack\!\rbrack}
\providecommand{\euro}{{\text{€}}}
\providecommand{\no}{{\text{n°}}}
\providecommand{\interval}[2]{{[#1;#2]}}
\providecommand{\intervFO}[2]{{[#1;#2[}}
\providecommand{\dotfill}{\ldots}

\providecommand{\ds}{\displaystyle}
\providecommand{\textding}[1]{#1}
\providecommand{\Oij}{\(\mathrm{O};\vec{\imath},\vec{\jmath}\)}
\providecommand{\Oijk}{\(\mathrm{O};\vec{\imath},\vec{\jmath},\vec{k}\)}
\providecommand{\Ouv}{\(\mathrm{O};\vec{u},\vec{v}\)}
\providecommand{\intervFF}[2]{\ensuremath{[#1;#2]}}
\renewcommand{\euro}{\ifmmode\text{\texteuro}\else\texteuro\fi}
\DeclareUnicodeCharacter{00B0}{\ifmmode^{\circ}\else\textdegree\fi}
\DeclareUnicodeCharacter{20AC}{\ifmmode\text{\texteuro}\else\texteuro\fi}
% Énoncés : Éducation nationale. Transcription et corrigés : APMEP, Denis Vergès et contributeurs. Découpage et classement : CoopMaths, Sébastien Lozano, Éric Elter et contributeurs. Adaptation de la présentation : Maths Explorer.
% CC BY-SA : https://coopmaths.fr/www/about/
% https://coopmaths.fr/alea/static/bac/2021/tex/bac_2021_01_sujet0_5.tex
\begin{document}
\noindent\textit{APMEP / CoopMaths. CC BY-SA. Présentation adaptée par Maths Explorer.}\par\medskip
\medskip

\framebox{
\begin{minipage}[]{5.25cm}
\textbf{Principaux domaines abordés}\\
\hspace{2em}Équations différentielles\\
\hspace{2em}Fonction exponentielle ; suites
\end{minipage}
}

\medskip

Dans une boulangerie, les baguettes sortent du four à une température de 225 \textcelsius.

On s’intéresse à l’évolution de la température d’une baguette après sa sortie du four.

On admet qu’on peut modéliser cette évolution à l’aide d’une fonction $f$ définie et dérivable sur l’intervalle $[0~;~+\infty[$.

 Dans cette modélisation, $f(t)$ représente la température en degré Celsius de la baguette au bout de la durée $t$, exprimée en heure, après la sortie du four.

Ainsi,$f(0,5)$ représente la température d’une baguette une demi-heure après la sortie du four.

Dans tout l’exercice, la température ambiante de la boulangerie est maintenue à 25 \textcelsius.

On admet alors que la fonction $f$ est solution de l'équation différentielle $y'+ 6y = 150$.

\medskip

\begin{enumerate}
\item 
	\begin{enumerate}
		\item Préciser la valeur de $f(0)$.
		\item Résoudre l’équation différentielle $y'+6y = 150$.
		\item En déduire que pour tout réel $t\geqslant  0$, on a $f(t) = 200 \e^{-6t}+25$.
	\end{enumerate}
\item Par expérience, on observe que la température d’une baguette sortant du four :

\begin{itemize}
\item  décroît ;
\item tend à se stabiliser à la température ambiante.
\end{itemize}

La fonction $f$ fournit-elle un modèle en accord avec ces observations ?

\item Montrer que l’équation $f(t) = 40$ admet une unique solution dans $[0~;~+\infty[ $.

Pour mettre les baguettes en rayon, le boulanger attend que leur température soit inférieure ou égale à 40 \textcelsius. On note $\mathcal{T}_0$ le temps d’attente minimal entre la sortie du four d’une baguette et sa mise en rayon.

On donne en page suivante la représentation graphique de la fonction $f $dans un repère orthogonal.

\begin{center}
\psset{xunit=5cm,yunit=0.025cm,labelFontSize=\scriptstyle,comma=true,labelsep=0.1pt}
\begin{pspicture}(-0.4,-20)(2.50,260)
\multido{\n=0.0+0.1}{23}{\psline[linewidth=0.3pt,linecolor=lightgray](\n,-20)(\n,260)}
\multido{\n=0+20}{14}{\psline[linewidth=0.3pt,linecolor=lightgray](0,\n)(2.1,\n)}
\psaxes[linewidth=0.95pt,Dx=0.5,Dy=20]{->}(0,0)(0,0)(2.1,260)
\psplot[linewidth=1.25pt,linecolor=blue,plotpoints=5000]{0}{2.1}{2.71828 x 6 mul neg   exp 200 mul 25 add}
\uput[d](1.75,-15){\footnotesize Durée en heure}
\uput[r](0,250){\footnotesize Température en degré Celsius}
\uput[ur](0.3,60){\blue $\mathcal{C}_f$}
\end{pspicture}
\end{center}

\item Avec la précision permise par le graphique, lire $\mathcal{T}_0$. On donnera une valeur approchée de $\mathcal{T}_0$ sous forme d’un nombre entier de minutes.

\item On s’intéresse ici à la diminution, minute après minute, de la température d’une baguette à sa sortie du four.

Ainsi, pour un entier naturel $n$, $\mathcal{D}_n$ désigne la diminution de température en degré Celsius d’une baguette entre la $n$-ième et la $(n+1)$-ième minute après sa sortie du four.

On admet que, pour tout entier naturel $n$ : 

\[\mathcal{D}_n=f\left(\dfrac{n}{60}\right)-f\left(\dfrac{n+1}{60}\right).\]

\smallskip

	\begin{enumerate}
		\item Vérifier que 19 est une valeur approchée de $\mathcal{D}_0$ à 0,1 près, et interpréter ce résultat dans le contexte de l’exercice.
\item Vérifier que l’on a, pour tout entier naturel $n$: 

\[\mathcal{D}_n = 200 \e^{-0,1n}\left(1 - \e^{-0,1}\right).\]

En déduire le sens de variation de la suite $\left(\mathcal{D}_n\right)$, puis la limite de la suite $\left(\mathcal{D}_n\right)$.

Ce résultat était-il prévisible dans le contexte de l’exercice ?

		\end{enumerate}
\end{enumerate}

\end{document}
